\section{Preliminaries}
\label{sec:prelim}

% We denote client indices as $j\in\{1,\dots,\Nc\}$ and class indices as $c\in\{1,\dots,\Cc\}$.

We denote client indices with $j\in\{1,\dots,\Nc\}$ and class indices are
$c\in\{1,\dots,\Cc\}$. Client $j$ holds a dataset $\Dj$ of size $\nj$ and uploads
a head displacement $\disp$, from which the server forms the shared head
$\thstar$.The frozen public backbone is $\featx$, while $\varphi_j$ denotes this backbone composed with client $j$'s local adapter. We denote $\mathsf{Enc}(\cdot)$
for encryption under the collective public key and $\tc$ for the decryption
threshold. \trsee{sec:prelim} provides a comprehensive table of the notation.

\tronly{\begin{table}[t]
\centering
\caption{Notation.}
\label{tab:notation}
\begin{tabular}{ll}
\toprule
Symbol & Meaning \\
\midrule
$\Nc$ & number of participating clients \\
$\Dj,\ \nj$ & client $j$'s dataset and its size $\nj=|\Dj|$ \\
$\Cc$ & number of classes \\
$n_{j,c}$ & examples of class $c$ held by client $j$ \\
% $\alpha$ & Dirichlet concentration, the label-skew parameter \\
$\featx$ & frozen public backbone, never trained or encrypted \\
$A_j$ & client $j$'s adapter, trained locally and never transmitted \\
$\varphi_j$ & client $j$'s feature map, the backbone composed with $A_j$ \\
$\thz$ & public initializer of the classifier head \\
$h_j$ & client $j$'s trained head \\
$\disp$ & head displacement $h_j-\thz$ \\
$\wj$ & sample weight $\nj/\sum_i n_i$ \\
$\thstar$ & shared head, held only as a ciphertext \\
$\mathsf{Enc}(\cdot)$ & encryption under the collective public key \\
\bottomrule
\end{tabular}
\end{table}}

\subsection{Multiparty Homomorphic Encryption}
\label{sec:mhe}

We build on CKKS~\cite{cheon2017ckks}, which performs approximate arithmetic on
packed vectors of real numbers under encryption. It is
a levelled scheme\tronly{~\cite{lyubashevsky2010ideal}} such that each multiplication spends one
level of a finite budget\tronly{, and a depleted budget is restored by
bootstrapping~\cite{cheon2018bootstrapping}}.

A single-key scheme would require one party to hold the secret key and therefore
to be trusted with decryption. We use the multiparty
variant~\cite{mouchet2021multiparty}, stated here for a $\tc$-out-of-$\Nc$
threshold with $2\le\tc\le\Nc$. 

The protocol requires four operations from the multiparty scheme. Key
generation produces one collective public key and the evaluation keys from the
clients' shares, and never assembles the secret in one place, such that any party may
compute on ciphertexts, and none may read them. Decryption takes partial shares
from at least $\tc$ clients, and any $\tc-1$ of them learn nothing about a
plaintext. A key switch to a designated public key redirects a result so that one
chosen party reads it and no other does. A collective refresh restores a spent
level budget, and enables bootstrapping. 

In the protocol, any $\tc$ clients form a
quorum, thus computation does not require every client to be online. Any $\tc-1$
clients are powerless against a ciphertext, hence $\tc-1$ is the corruption bound
throughout \cref{sec:security}. 

Our implementation instantiates $\tc=\Nc$, which is the setting the multiparty
CKKS library provides, and the measurements of \cref{sec:exp-cost} are taken
there. Every statement in \cref{sec:security} holds for general $\tc$.

%We state the reason for this choice. The aggregation of \cref{sec:split} uses only additions and multiplications by public scalars, hence it consumes one level. What we need from the scheme is not depth but key structure: threshold decryption over a collectively generated key is the property the protocol cannot do without, and mature distributed key generation and collective key switching provide it directly. CKKS additionally packs thousands of real values into one ciphertext and multiplies them by public scalars natively. Additively homomorphic schemes offer neither packed real arithmetic at this width nor comparable multiparty tooling, and masking-based secure aggregation reveals the plaintext sum to the server, which this protocol never does.
